\scanpage{131}
% archive sheet number: 130

Cette situation permet de définir
sur $P$

\[
\begin{aligned}
&\text{forme linéaire }L_P=\mathrm{Tr}_A\circ f,\ \text{qui prolonge }T_B\\
&\text{forme quadratique }\Delta_P=\det_A\circ f\\
&\text{forme bilinéaire sym. }\psi_P\ \text{polarisée de }\Delta_P,\ \psi_P=\psi_A\circ f\\
&\hspace{6em}\text{(où }\psi_A\text{ est polarisée de }\det_A\text{)}\\
&\text{forme bilinéaire }\varphi_P=\varphi_A\circ f
\end{aligned}
\]

i.e.

\[
\varphi_P\deleted{\unsure{=}}(U,V)=\mathrm{Tr}\bigl(f(U)\unsure{,|.}f(V)\bigr)
\]

\margin{\uline{NB} On a \deleted{$\unsure{\Delta_P,\varphi_P,\psi_P}$} $\varphi_P(U,V)+\psi_P(U,V)=L_P(u)L_P(V)$ et $\varphi_P(1,U)=\psi_P(1,U)=L_P(U)$.}

D'autre part, on peut restreindre ceci
à $B$, on trouve

\[
\text{forme quad }\Delta_B,\quad\text{formes bil. sym }\psi_B,\varphi_B.
\]

En restreignant ces trois formes sur
$B$ à $H$, on trouve les trois formes
dont $H$ est déjà équipé, $\Delta_H=\det_H$, $\psi_H=$ forme
polarisée de $\Delta_H$, $\varphi_H=\bigl((u,v)\unsure{-|\mapsto}\mathrm{Tr}_H(uv)\bigr)$, Et
itou pour $H'$.

On a d'autre part

\[
\Delta_P(u\otimes v)=\Delta_H(u)\,\Delta_{H'}(v) \tag*{(1)}
\]

\uline{Proposition} \inserted{On se donne $H,H'$ comme dessus.} Considérons les \deleted{\unsure{trois|six}} \inserted{quatre} ensembles
suivants :

$Q_P$: \deleted{$(\unsure{\Delta_P})$} ens des formes quadratiques $\Delta_P$ sur $P$ satisfaisant
(1) ci-dessus

$Q_B$: \deleted{$(\unsure{\Delta_B|\Delta_H})$} ens des formes \deleted{\unsure{bil.}} \unsure{qdr.|quadr.} $\Delta_B$ sur $B$ qui
prolongent $\Delta_H$ et $\Delta_{H'}$

$\underline{\Psi}_B$: \deleted{$(\unsure{\psi})$} ens des formes bil. sym $\psi_B/\Delta_B$ sur $B$
qui prolongent $\psi_H,\psi_{H'}$

$\underline{\Phi}_B$: \deleted{$(\varphi)$} ens des formes bil. symétriques $\varphi_B$ sur $B$

